The vacuum term in a collision contributes to chiral chains on the projective line.
For the critical Virasoro algebra with ghosts, its differential admits a simpler description through global vector fields.
The resulting cohomology has generators in degrees one and three.
Both the collision differential and the degree shifts enter this calculation.

\section{The critical differential and the ghost subalgebra}

Work over $\mathbb C$, with algebraic tensor products and direct sums.
Set $X=\mathbb P^1$.
Complexes are cohomological, with $K[r]^j=K^{j+r}$ and differential $(-1)^r d_K$.
The sign for interchanging homogeneous elements of degrees $a,b$ is $(-1)^{ab}$.
Ghost degree gives both cohomological degree and parity modulo two.
Conformal weight is a separate integer grading.

Let $V=\operatorname{Vir}_{26}$ be the universal Virasoro vacuum vertex algebra.
Its vacuum $\mathbf1$ satisfies $L_n\mathbf1=0$ for $n\geq-1$.
The ordered monomials in $L_{-n}$, $n\geq2$, form its vector-space basis.
Write $T=L_{-2}\mathbf1$.
Let $\mathcal F$ be the $bc$ vertex superalgebra with
\[
 \{b_m,c_n\}=\delta_{m+n,0},\qquad
 \{b_m,b_n\}=\{c_m,c_n\}=0,
 \qquad
 b_n\Omega=0\ (n\geq-1),\quad c_n\Omega=0\ (n\geq2).
\]
Its creators are $b_n$ for $n\leq-2$ and $c_n$ for $n\leq1$.
The generating states are $b=b_{-2}\Omega$ and $c=c_1\Omega$, of degrees $-1,1$ and weights $2,-1$.
Put $A=V\otimes\mathcal F$ and write $\Omega$ also for its vacuum.
The field conventions are
\begin{gather*}
 Y(T,z)=\sum_nL_nz^{-n-2},\qquad
 Y(b,z)=\sum_nb_nz^{-n-2},\\
 Y(c,z)=\sum_nc_nz^{1-n},\qquad
 Y(a,z)=\sum_na_{(n)}z^{-n-1}.
\end{gather*}
In particular, the vertex product $b_{(0)}$ is the conformal mode $b_{-1}$ and differs from the constraint operator $b_0$.

Normal ordering moves all creators to the left with their odd permutation signs and omits contractions.
On $A$, define
\begin{equation}\label{eq:uc-Q}
 Q=\sum_n c_{-n}L_n
 -\frac12\sum_{m,n}(m-n):c_{-m}c_{-n}b_{m+n}:.
\end{equation}
These sums are locally finite on each algebraic state.
The total stress tensor and its zero mode are
\[
 T^{\mathrm{tot}}=T-2:b\partial c:-:\partial b\,c:,
 \qquad H=L_0+L_0^{\mathrm{gh}}.
\]
Write $Y(T^{\mathrm{tot}},z)=\sum_n\mathcal L_nz^{-n-2}$, so $H=\mathcal L_0$.
Theorem~\ref{thm:vir-square} gives $Q^2=0$ and $\{Q,b_0\}=H$.
The field in~\eqref{eq:vir-charge-residue} has zero mode $Q$.
The vertex commutator identity for a zero mode therefore gives
\begin{equation}\label{eq:uc-derivation}
 Q(a_{(n)}a')=(Qa)_{(n)}a'+(-1)^{|a|}a_{(n)}Qa'.
\end{equation}
In particular, $Q$ commutes with translation $\partial$, and
\begin{equation}\label{eq:uc-Qfields}
 Qc=:c\partial c:,
 \qquad Qb=T^{\mathrm{tot}},
 \qquad QT=:c\partial T:+2:\partial c\,T:+\frac{13}{6}\partial^3c.
\end{equation}
For the last identity, take the residue of $c(z)T(z)T(w)$.
The coefficients of $(z-w)^{-4},(z-w)^{-2},(z-w)^{-1}$ are $13,2T,\partial T$.
Taylor expansion of $c(z)$ gives the stated coefficients.

Let $B\subset A$ be the subalgebra generated by $c$ and its derivatives.
It is the commutative vertex algebra
\begin{equation}\label{eq:uc-B}
 B=\Lambda(c^{(j)}:j\geq0),\qquad
 c^{(j)}=\partial^jc=j!c_{1-j}\Omega,
 \qquad Qc^{(j)}=\partial^j(c\partial c).
\end{equation}
Every generator has degree one.
All singular products between these generators vanish.
The regular vertex operation is $Y(f,t)g=(e^{t\partial}f)g$.

The total central charge is zero.
The identities $[Q,\mathcal L_n]=0$ in~\eqref{eq:vir-Q-L} permit coordinate descent of $(A,Q)$.
The $c$-ghosts are preserved by these coordinate changes.
For a change $w=w(z)$ their field transforms by
\begin{equation}\label{eq:uc-c-transition}
 c_w(w)=\frac{dw}{dz}c_z(z).
\end{equation}
Repeated differentiation determines every jet transition in $B$.
Denote the resulting left state modules by $\mathcal A^\ell,\mathcal B^\ell$.
The second is a commutative differential graded $\mathcal D_X$-algebra.
Their right modules are $\mathcal A=\omega_X\otimes\mathcal A^\ell$ and $\mathcal B=\omega_X\otimes\mathcal B^\ell$.
They carry the usual chiral products associated with their vertex algebras.
In a coordinate, the connection on a state section is $\partial_z+\partial$.
Thus the line generated by $c$ in the left module is $\omega_X$.
The connection and the primary-state line convention agree with Frenkel~\cite{Frenkel026}, \S\S4.2--4.3, pp.~875-18--875-19.
The conformal weights here are bounded below by $-1$.
Positive Virasoro modes lower weight and hence act locally nilpotently, so the same coordinate-bundle construction applies.

\begin{proposition}\label{prop:uc-B-qiso}
The inclusion $\mathcal B\longrightarrow\mathcal A$ is a quasi-isomorphism of unital differential graded chiral superalgebras.
\end{proposition}

\begin{proof}
The formulas above prove compatibility with the unit, products, differential, and coordinate changes.
It remains to calculate the cohomology in a coordinate trivialization.
Both $A$ and $B$ have integral energy decompositions with finite energy support on each element.
On a nonzero energy $\lambda$, the operator $\lambda^{-1}b_0$ contracts both complexes.
On $B$, $b_0$ differentiates the exterior variable $c_0$ and preserves $B$.
The contraction identity is the restriction of $\{Q,b_0\}=H$.

The matter algebra has $V_0=\mathbb C\mathbf1$, $V_1=0$, and no negative weights.
The only negative-energy ghost creator is $c_1$, which occurs at most once.
Every $b$ creator has energy at least two.
Consequently the energy-zero subspaces of $A$ and $B$ are identical, with basis
\[
 e=\Omega,\qquad u=c_0\Omega,\qquad
 v=c_{-1}c_1\Omega,\qquad w=c_{-1}c_0c_1\Omega.
\]
Their degrees are $0,1,2,3$, and the differential is
\begin{equation}\label{eq:uc-zeroenergy}
 Qe=0,\qquad Qu=-2v,\qquad Qv=Qw=0.
\end{equation}
Here $Qu$ follows from $\{Q,c_0\}=-2\sum_{n>0}n c_{-n}c_n$.
Nilpotence gives $Qv=0$, and degree gives $Qw=0$.
The map $h(v)=-u/2$, zero on the other three basis vectors, contracts the pair $u,v$.
Thus both cohomologies have basis $[e],[w]$, and inclusion identifies them.
This local calculation proves the assertion for the underlying sheaf complexes.
The contraction need not be invariant under coordinate changes for this purpose.
\end{proof}

\section{Global vector fields and a constant differential graded algebra}

Let $\mathfrak g=H^0(X,T_X)$ with its vector-field bracket.
On the affine chart with coordinate $z$, use
\[
 e_- =\partial_z,\qquad e_0=z\partial_z,\qquad e_+=z^2\partial_z.
\]
Then $[e_-,e_0]=e_-$, $[e_-,e_+]=2e_0$, and $[e_0,e_+]=e_+$.
Let $E=C^\bullet_{\mathrm{CE}}(\mathfrak g,\mathbb C)$ be its Chevalley--Eilenberg cochain algebra.
If $\alpha,\beta,\gamma$ are the dual degree-one generators, its differential is
\begin{equation}\label{eq:uc-E}
 E=\Lambda(\alpha,\beta,\gamma),\qquad
 d\alpha=-\alpha\beta,\quad d\beta=-2\alpha\gamma,\quad
 d\gamma=-\beta\gamma.
\end{equation}
Write $E_X^\ell=\mathcal O_X\otimes E$ with the connection on the first factor, and $E_X=\omega_X\otimes E$.

\begin{proposition}\label{prop:uc-evaluation}
There is a quasi-isomorphism of unital differential graded chiral algebras
\[
 \epsilon:\mathcal B\longrightarrow E_X.
\]
In the coordinate $z$, it is determined on left modules by
\begin{equation}\label{eq:uc-eval}
 \epsilon(c)=-f(z),\qquad
 f(z)=\alpha+z\beta+z^2\gamma,
 \qquad \epsilon(\partial^jc)=\partial_z^j\epsilon(c).
\end{equation}
\end{proposition}

\begin{proof}
The last formula is the equation $\partial_z\epsilon(c)=\epsilon(\nabla_{\partial_z}c)$ for a morphism of left $\mathcal D_X$-modules.
Thus $\epsilon(\partial c)=-f'(z)$ and $\epsilon(\partial^2c)=-2\gamma$.
Equation~\eqref{eq:uc-E} gives $df=-ff'$.
Consequently
\[
 d\epsilon(c)=ff',\qquad
 \epsilon(Qc)=(-f)(-f')=ff'.
\]
The two differentials commute with translation, so they agree on every jet.
Both are derivations, so they agree on $B$.

The vector field $f(z)\partial_z$ is the tautological element of $\mathfrak g\otimes\mathfrak g^*$, with its coefficients placed in degree one.
Under $w=w(z)$, its coefficient becomes $f_w=w'f_z$.
This agrees with~\eqref{eq:uc-c-transition}, and the jet equations agree by differentiation.
For the two standard charts $w=1/z$, one obtains
\[
 \epsilon(c_w)=\alpha w^2+\beta w+\gamma
 =(-z^{-2})\epsilon(c_z).
\]
The formulas are regular in both charts.
They therefore define the claimed global morphism.
Commutative vertex products are determined by multiplication and translation, so this is also a chiral-algebra morphism.

The differential $E^1\to E^2$ in~\eqref{eq:uc-E} is an isomorphism.
Its matrix in the ordered bases $(\alpha,\beta,\gamma)$ and $(\alpha\beta,\alpha\gamma,\beta\gamma)$ is $\operatorname{diag}(-1,-2,-1)$.
The other nonzero cohomology classes of $E$ are $1$ and $\eta=\alpha\beta\gamma$.
At $z=0$,
\[
 \epsilon(w)=\epsilon\bigl(\tfrac12(\partial^2c)(\partial c)c\bigr)
 =(-\gamma)(-\beta)(-\alpha)=\eta.
\]
The same product equals $\eta$ at every finite $z$.
In the $w$ chart its value is also $\eta$.
The four-state calculation in~\eqref{eq:uc-zeroenergy} therefore proves that $\epsilon$ induces an isomorphism on every cohomology stalk.
\end{proof}

Let $K=\Lambda(\eta)$ with $|\eta|=3$ and zero differential.
The map $K\to E$ given by $\eta\mapsto\alpha\beta\gamma$ is a unital quasi-isomorphism of commutative differential graded algebras.
Combining the two propositions gives a zigzag of actual chiral-algebra morphisms
\begin{equation}\label{eq:uc-zigzag}
 \mathcal A\ \longleftarrow\ \mathcal B\
 \xrightarrow{\epsilon}\ E_X\ \longleftarrow\ K_X,
 \qquad K_X=\omega_X\otimes K.
\end{equation}
Every arrow is a quasi-isomorphism.
The arrows have their displayed directions: a chosen inverse exists only after passing to the derived category.

\section{The Chevalley--Cousin differential}

For a nonempty finite set $I$, let $\operatorname{Part}(I)$ be its set of partitions.
Regard a partition $\pi$ as the quotient map from $I$ to its set of blocks.
Let $\Delta_\pi:X^\pi\hookrightarrow X^I$ identify coordinates within each block.
Let $U^\pi\subset X^\pi$ be the open set of distinct block coordinates, with inclusion $j_\pi$.
For a right differential graded chiral algebra $\mathcal M$, define
\begin{equation}\label{eq:uc-Cousin}
 C(\mathcal M)_{X^I}
 =\bigoplus_{\pi\in\operatorname{Part}(I)}
 (\Delta_\pi)_*(j_\pi)_*j_\pi^*
 (\mathcal M[1])^{\boxtimes\pi}.
\end{equation}
Here $\Delta_*$ is direct image for right $\mathcal D$-modules.
The open direct image allows poles of arbitrary finite order along the omitted diagonals.
Every summand is algebraic: no inverse limit in pole order or arity is taken.
Formula~\eqref{eq:uc-Cousin} is Beilinson--Drinfeld's complex~\cite{BD026}, equation (3.4.11.1), p.~182.

Write $sa$ for an element of $\mathcal M[1]$, so $|sa|=|a|-1$.
The internal differential is the tensor differential extending $sa\mapsto-s(d_\mathcal M a)$.
The collision differential merges two blocks using the suspended chiral product
\begin{equation}\label{eq:uc-suspended}
 d_2(sa\,sb)=(-1)^{|a|}s\mu(a,b).
\end{equation}
For any pair of blocks, move their suspended factors to the first two positions using the Koszul rule, apply~\eqref{eq:uc-suspended}, and retain the other factors.
The sum over all pairs defines $d_{\mathrm{col}}$.
In~\eqref{eq:uc-suspended}, the argument is localized away from the two-block diagonal and the output is its direct image.
Localizations at the other diagonals remain in place.
This specifies a morphism of right $\mathcal D_{X^I}$-modules on each summand.

\begin{proposition}\label{prop:uc-Ccomplex}
With $d=d_{\mathrm{int}}+d_{\mathrm{col}}$, formula~\eqref{eq:uc-Cousin} defines a complex.
The diagonal structure maps and the factorization identifications commute with $d$.
The construction is functorial for differential graded chiral-algebra morphisms.
\end{proposition}

\begin{proof}
The internal square is zero by the differential on $\mathcal M$.
Equation~\eqref{eq:uc-derivation}, with the suspension signs in~\eqref{eq:uc-suspended}, gives
$d_{\mathrm{int}}d_{\mathrm{col}}+d_{\mathrm{col}}d_{\mathrm{int}}=0$.
Two mergers of disjoint pairs cancel in opposite orders because each merger has degree one.
For three blocks, the three terms in the remaining square are the suspended chiral Jacobi identity.
These cases exhaust all terms in $d_{\mathrm{col}}^2$.

A diagonal structure map includes the summands whose partitions factor through the specified quotient of labels.
Merging blocks commutes with this inclusion.
Over separated groups of points, partitions split into the partitions of those groups.
The tensor-product identification of the corresponding summands commutes with internal differentials and collisions inside each group.
For three groups, both composite identifications use the same partition and the same tensor symmetry.
They therefore agree.
The same argument with a morphism on every block proves functoriality.
\end{proof}

For $|I|=2$, the complex contains precisely
\[
 j_*j^*(\mathcal M[1]\boxtimes\mathcal M[1])
 \xrightarrow{d_{\mathrm{col}}}\Delta_*\mathcal M[1].
\]
For $|I|=3$, its terms have three, two, and one blocks.
The two-block term has one summand for each pair of labels.
Thus the first coherence after a binary collision is the three-block Jacobi identity, with all three intermediate diagonals present.

\section{An actual global chain complex}

The de Rham functor on right $\mathcal D$-modules has the convention
\[
 \operatorname{DR}_X(\mathcal M)
 =[\mathcal M\otimes_{\mathcal O_X}T_X\longrightarrow\mathcal M]
\]
in auxiliary degrees $-1,0$.
In particular, $\operatorname{DR}_X(\omega_X)=\Omega_X^\bullet[1]$.
To keep collision operations strict after taking sections, choose the following two resolutions.

Let $P^0=\operatorname{Sym}_{\mathcal O_X}^{>0}(\mathcal D_X)$, where the left $\mathcal D_X$-action is the diagonal action.
Its augmentation sends a differential operator $D$ in the generating summand to $D(1)$.
Set
\[
 P^{-1}=\ker(P^0\to\mathcal O_X),\qquad d_P:P^{-1}\hookrightarrow P^0.
\]
Multiplication in degree zero is the symmetric-algebra multiplication.
Its action on $P^{-1}$ is the ideal action, and the product of two degree-minus-one elements is zero.
The Leibniz identity on two such elements is $uv-uv=0$.
This defines a nonunital commutative differential graded $\mathcal D_X$-algebra resolution of $\mathcal O_X$.
Its terms are $\mathcal D_X$-flat by the curve construction in~\cite{BD026}, \S2.2.10, p.~72.
Nonunitality belongs to this auxiliary resolution and imposes no augmentation on $\mathcal A$.

Use the two affine charts $U_0,U_\infty$ and their intersection $U_{01}$.
Write $j_i,j_{01}$ for their embeddings in $X$.
Let $\mathcal Q^0$ consist of triples
\[
 (f_0,f_\infty,F(s))\in
 j_{0*}\mathcal O_{U_0}\oplus j_{\infty*}\mathcal O_{U_\infty}
 \oplus j_{01*}\mathcal O_{U_{01}}[s]
\]
such that $F(0)=f_0|_{U_{01}}$ and $F(1)=f_\infty|_{U_{01}}$.
Set $\mathcal Q^1=j_{01*}\mathcal O_{U_{01}}[s]\,ds$, with differential $F\mapsto F'(s)ds$.
Multiplication is polynomial multiplication with $ds^2=0$.
The connection differentiates the coefficients, holding $s$ fixed.
This is the two-chart Thom--Sullivan resolution from~\cite{BD026}, \S4.1.3, pp.~278--279.
Each term is $\mathcal O_X$-flat.
For $\mathcal Q^0$, this also follows from its exact sequence with quotient
$j_{0*}\mathcal O_{U_0}\oplus j_{\infty*}\mathcal O_{U_\infty}$ and kernel
$s(s-1)j_{01*}\mathcal O_{U_{01}}[s]$.

For any $\mathcal M$, put $\mathcal M_{P\mathcal Q}=\mathcal M\otimes P\otimes \mathcal Q$.
The tensor differential on $a\otimes p\otimes q$ is
\[
 d_\mathcal M a\otimes p\otimes q
 +(-1)^{|a|}a\otimes d_Pp\otimes q
 +(-1)^{|a|+|p|}a\otimes p\otimes d_{\mathcal Q}q.
\]
Tensor symmetry and multiplication in $P,\mathcal Q$ extend the chiral product.
For a right $\mathcal D_Y$-module $N$, set $h_Y(N)=N\otimes_{\mathcal D_Y}\mathcal O_Y$, applied degreewise.
Define
\begin{equation}\label{eq:uc-chain}
 \mathsf C_X(\mathcal M)
 =\underset{I}{\operatorname{colim}}
 \Gamma\bigl(X^I,h_{X^I}(C(\mathcal M_{P\mathcal Q})_{X^I})\bigr).
\end{equation}
The index category has nonempty finite sets and arrows opposite to surjections.
Its transition maps are the diagonal structure maps.
Thus~\eqref{eq:uc-chain} is an actual complex of vector spaces with the differential already specified.
The comparison of Beilinson--Drinfeld, equation (4.2.12.2), pp.~305--306, identifies its derived-category object with
\[
 C^{\mathrm{ch}}(X,\mathcal M)=R\Gamma_{\mathrm{DR}}(\operatorname{Ran}(X),C(\mathcal M)).
\]
Their comparison is functorial and preserves the filtration by the number of blocks.

As a graded vector space,~\eqref{eq:uc-chain} has the smaller description
\begin{equation}\label{eq:uc-small}
 \mathsf C_X(\mathcal M)
 =\bigoplus_{n\geq1}
 \Gamma\bigl(U^{[n]},h_{X^n}((\mathcal M_{P\mathcal Q}[1])^{\boxtimes n})\bigr)_{\mathfrak S_n}.
\end{equation}
This is equation (4.2.12.5), p.~307, with its induced differential.
The symmetric group acts on the suspended factors.
The subscript means coinvariants over $\mathbb C$.
The summand with one point contains vacuum states through the unit $\omega_X\to\mathcal M$.
There is no arity-zero summand in this reduced Chevalley complex.

\section{The first collision and its differential}

Fix a formal coordinate $t=z_1-z_2$ transverse to the diagonal.
Normalize $\operatorname{Res}(t^{-1}dt)=1$.
For a Laurent polynomial $f$ and states $a,a'$, the normal residue coefficient of the collision is
\begin{equation}\label{eq:uc-rho}
 \rho_f(a,a')=\operatorname{Res}_{t=0}f(t)Y(a,t)a'\,dt.
\end{equation}
The whole distribution is its principal part $[f(t)Y(a,t)a']$.
This is the local chiral-product formula of~\cite{Frenkel026}, Theorem~7.1 and Corollary~7.2, p.~875-35.
Pairing this principal part with all powers of $t$ recovers every normal derivative.
Lower truncation makes each principal part finite.
The vacuum and Virasoro products give
\begin{equation}\label{eq:uc-firstcoll}
 \rho_1(b,c)=1,\qquad \rho_{t^3}(T,T)=13\,1,
 \qquad \rho_{t^{-1}}(1,a)=a.
\end{equation}
All three identities occur in~\eqref{eq:uc-Cousin} and~\eqref{eq:uc-chain}.

The first mixed differential identity can be computed on $b,c$ without infinite sums:
\begin{align*}
 Q\rho_1(b,c)
 &=\rho_1(Qb,c)-\rho_1(b,Qc)\\
 &=(T^{\mathrm{tot}})_{(0)}c-b_{(0)}(:c\partial c:)\\
 &=\partial c-\partial c=0.
\end{align*}
The second term uses $b_{(0)}c=1$ and $b_{(0)}\partial c=0$.
The first uses the translation mode of the total stress tensor.
For the Virasoro collision,
\[
 (QT)_{(3)}T=-26\partial c,
 \qquad T_{(3)}QT=26\partial c.
\]
Indeed, differentiating the central term $13t^{-4}$ in the first variable gives $-52t^{-5}$.
Its linear Taylor coefficient contributes $-52\partial c$.
The term $2\partial c\,T$ contributes $26\partial c$.
In the second expression, only $2\partial c\,T$ contributes to the coefficient of $t^{-4}$.
This proves $Q\rho_{t^3}(T,T)=0$ directly.

The general identity
\[
 Q\rho_f(a,a')=\rho_f(Qa,a')+(-1)^{|a|}\rho_f(a,Qa')
\]
follows coefficientwise from~\eqref{eq:uc-derivation}.
It also proves binary compatibility for every arrow of~\eqref{eq:uc-zigzag}.
For $\epsilon$, its compatibility with the regular collision $f=t^{-1}$ reads
$\epsilon(aa')=\epsilon(a)\epsilon(a')$.
Compatibility with all other $f$ follows from its translation equation.
The first three-point compatibility is therefore the commuting diagram obtained by applying each arrow to all three Jacobi composites.

\section{The filtration and its chiral symbol}

Use the Poincar\'e--Birkhoff--Witt (PBW) ordering on the negative matter modes and the exterior ghost generators.
Give each negative matter mode and each $b$ creator order one.
Give every $c$ creator order zero.
Let $F_pA$ be the span of ordered monomials of order at most $p$, with $F_{-1}A=0$.
This is an increasing, exhaustive filtration, with $F_0A=B$.
Translation and coordinate changes preserve it.
Normally ordered products add orders, and each singular contraction lowers their sum.
Equations~\eqref{eq:uc-Qfields} show that $Q$ preserves $F$.
Its action on an arbitrary ordered monomial follows from~\eqref{eq:uc-derivation} with $n=-1$.

The associated graded vertex algebra is commutative, with local description
\begin{equation}\label{eq:uc-symbolalg}
 R=\mathbb C[t^{(j)}:j\geq0]\otimes
 \Lambda(b^{(j)},c^{(j)}:j\geq0).
\end{equation}
Here $t,b,c$ are the symbols of $T,b,c$, their degrees are $0,-1,1$, and their orders are $1,1,0$.
The index $(j)$ means $j$ translations.
The differential $q$ commutes with translation and is determined by
\begin{equation}\label{eq:uc-symbolq}
 qc=c\partial c,\qquad
 qb=t-2b\partial c-(\partial b)c,\qquad
 qt=c\partial t+2(\partial c)t.
\end{equation}
The term $13\partial^3c/6$ in $QT$ has lower order and hence disappears from $qt$.
The singular collisions in~\eqref{eq:uc-firstcoll} also lower order.
Regular collisions remain: for example $\rho_{t^{-1}}(1,a)=a$ on $R$.

Put $\ell=qb$.
The change from $t$ to $\ell=t-2b\partial c-(\partial b)c$ is an invertible polynomial change of jet generators.
Equations~\eqref{eq:uc-symbolq} give $q\ell=0$ by direct substitution.
Thus
\[
 q= q_B+\sum_{j\geq0}\ell^{(j)}\frac{\partial}{\partial b^{(j)}}
\]
in these generators.
Let $N$ count the $\ell$ and $b$ factors and define the degree-minus-one derivation
\[
 h=\sum_{j\geq0}b^{(j)}\frac{\partial}{\partial\ell^{(j)}}.
\]
Each sum is finite on a polynomial, and $qh+hq=N$.
On the positive-order subspace, $N^{-1}h$ is a contraction.
Both $b$ and $\ell=qb$ transform as primary fields of weight two, so this construction respects the jet transition maps.
It follows that $\mathcal B\hookrightarrow\operatorname{gr}_F\mathcal A$ is a quasi-isomorphism, with positive orders acyclic.

Give~\eqref{eq:uc-Cousin} and~\eqref{eq:uc-chain} total order $\sum_i p_i$ on a tensor of factors in $F_{p_i}\mathcal A$.
Keep $P,\mathcal Q$ in order zero.
The symbol operation on each tensor induces the typed isomorphism of actual complexes
\begin{equation}\label{eq:uc-symbolmap}
 \sigma:\operatorname{gr}_F\mathsf C_X(\mathcal A)
 \xrightarrow{\ \cong\ }
 \mathsf C_X(\operatorname{gr}_F\mathcal A),
 \qquad
 [\boxtimes_i(a_i\otimes p_i\otimes q_i)]
 \longmapsto\boxtimes_i([a_i]\otimes p_i\otimes q_i).
\end{equation}
This is the chain-level comparison in~\cite{BD026}, equation (4.2.18.3), p.~311.
Its hypotheses hold because the chosen $P,\mathcal Q$ have flat terms and the state filtration has locally free associated graded pieces.
Localization, external tensor product, and diagonal direct image respect these symbols.
For the functor $h$ and global sections, the flat-resolution argument of \S4.2.12 gives the same comparison.
The internal differential becomes~\eqref{eq:uc-symbolq}.
The collision differential becomes the regular commutative vertex collision on $R$.
These identifications verify that~\eqref{eq:uc-symbolmap} commutes with both parts of the differential.

In particular, the first PBW target is the actual complex
\begin{equation}\label{eq:uc-grzero}
 \operatorname{gr}^F_0\mathsf C_X(\mathcal A)
 =\mathsf C_X(\mathcal B)
 =\bigoplus_{n\geq1}
 \Gamma\bigl(U^{[n]},h((\mathcal B_{P\mathcal Q}[1])^{\boxtimes n})\bigr)_{\mathfrak S_n},
\end{equation}
with~\eqref{eq:uc-B}, the $P,\mathcal Q$ differentials, and the regular collision products specified above.
For the separate filtration $G$ by the number of points,
\[
 \operatorname{gr}^G_1\operatorname{gr}^F_0\mathsf C_X(\mathcal A)
 =\Gamma(X,h(\mathcal B_{P\mathcal Q}[1])).
\]
Thus PBW order and collision length give different filtrations with different first terms.

There is also an actual filtered chain map
\begin{equation}\label{eq:uc-filtered-inclusion}
 \iota:\mathsf C_X(\mathcal B)\longrightarrow\mathsf C_X(\mathcal A),
\end{equation}
obtained by including each ghost factor in each partition summand.
Give its source the filtration concentrated in order zero.
Its associated graded map, followed by $\sigma$, is the map induced by $\mathcal B\hookrightarrow R$.
The contraction above and functoriality of chiral chains show that this associated graded map is a quasi-isomorphism.
Thus~\eqref{eq:uc-filtered-inclusion} is a filtered quasi-isomorphism.
This supplies a lift of the order-zero symbol on the full configuration complex.
The differential equation holds on every partition, including all binary and triple collisions.

\section{The cohomology on the projective line}

\begin{theorem}\label{thm:uc-global}
For $\mathcal A$ associated with $(\operatorname{Vir}_{26}\otimes bc,Q)$ on $X=\mathbb P^1$, the algebraic chiral chain complex satisfies
\begin{equation}\label{eq:uc-globalanswer}
 C^{\mathrm{ch}}(X,\mathcal A)\simeq\Lambda(\xi_1,\xi_3),
 \qquad |\xi_1|=1,\quad |\xi_3|=3,
\end{equation}
as an object of the derived category of complexes.
The algebra on the right has zero differential.
Its cohomology basis is $1,\xi_1,\xi_3,\xi_1\xi_3$, in degrees $0,1,3,4$.
With Beilinson--Drinfeld's homological convention $H_a^{\mathrm{ch}}=H^{-a}C^{\mathrm{ch}}$, the corresponding indices are $0,-1,-3,-4$.
The PBW associated graded chiral complex has the same cohomology, entirely in order zero.
\end{theorem}

\begin{proof}
Apply the functor $\mathsf C_X$ to~\eqref{eq:uc-zigzag}.
The result consists of actual chain maps, each a quasi-isomorphism by Beilinson--Drinfeld, \S4.2.11 and \S4.2.12.
It therefore suffices to compute $C^{\mathrm{ch}}(X,K_X)$.
The left algebra $K_X^\ell=\mathcal O_X\otimes\Lambda(\eta)$ is freely commutative on $\mathcal O_X[-3]$.
The corresponding right generator module is
\[
 W=\omega_X[-3].
\]
It is $\mathcal O_X$-flat.
Set $L=\Gamma(X,h(W_{P\mathcal Q}))[1]$.
The actual comparison is the following map of complexes:
\begin{equation}\label{eq:uc-free-map}
 \Psi:\mathsf C_X(\omega_X)\oplus\operatorname{Sym}^{>0}(L)
 \longrightarrow\mathsf C_X(K_X).
\end{equation}
On the first summand, use the chain map induced by the chiral unit $\omega_X\to K_X$.
On $L$, use the inclusion of the generator module into the one-point term.
For each positive symmetric degree, take the external product of these inclusions and restrict to the distinct-point locus.
The differential commutes with these products because $K_X$ is commutative.
This specifies $\Psi$ on every chain, including its entire first summand.

The positive-degree map factors through $\mathsf C_X(W)$, where $W$ has zero chiral product.
Indeed, $W$ is the nonunital free commutative chiral algebra on its odd rank-one generator, and $K_X=W\oplus\omega_X$ is its unitalization.
Equation (4.6.2.2) and the proof of the proposition in~\cite{BD026}, pp.~342--344, give a quasi-isomorphism
$\operatorname{Sym}^{>0}(L)\to\mathsf C_X(W)$.
The unit and inclusion induce another actual quasi-isomorphism of complexes
\[
 \mathsf C_X(\omega_X)\oplus\mathsf C_X(W)
 \longrightarrow\mathsf C_X(K_X)
\]
by the proof of~\cite{BD026}, \S4.4.8, p.~329.
Their composition is $\Psi$, which is therefore a quasi-isomorphism.

The unit contributes through the first complex summand.
Its replacement by $\mathbb C$ occurs in the derived category, using
$C^{\mathrm{ch}}(X,\omega_X)\simeq\mathbb C$ from~\cite{BD026}, \S4.3.3, p.~315.
The raw positive-arity chain multiplication has the homotopy-unit structure of \S4.3.4.
No representative of a cohomology class is required to define~\eqref{eq:uc-free-map}.

The complex $L$ represents
\[
 R\Gamma_{\mathrm{DR}}(X,W[1])
 =R\Gamma(X,\Omega_X^\bullet)[-1].
\]
Algebraic de Rham cohomology of $\mathbb P^1$ is $\mathbb C$ in degrees zero and two and vanishes otherwise.
For completeness, the two-chart \v Cech complex gives
$H^0(\mathcal O_X)=\mathbb C$, $H^1(\mathcal O_X)=0$,
$H^0(\Omega_X^1)=0$, and $H^1(\Omega_X^1)=\mathbb C[dz/z]$.
These are the only two entries in its de Rham spectral sequence, so no differential can join them.
Normalize $[dz/z]$ by the residue $1$ on the overlap.
After the displayed shift, these two classes have degrees one and three.
Over $\mathbb C$, symmetric powers preserve quasi-isomorphisms of complexes, because the symmetric-group coinvariant functors are exact.
Together with the first summand of~\eqref{eq:uc-free-map}, their positive symmetric powers give the exterior algebra in~\eqref{eq:uc-globalanswer}.

Finally,~\eqref{eq:uc-symbolmap} identifies the PBW symbol complex with the chiral chains of $R$.
The quasi-isomorphism $\mathcal B\to R$ preserves PBW order and has image in order zero.
Applying chiral chains and then~\eqref{eq:uc-zigzag} computes its cohomology by the same exterior algebra.
Every positive-order cohomology group therefore vanishes.
\end{proof}

The chiral-algebra zigzag~\eqref{eq:uc-zigzag}, the chain map~\eqref{eq:uc-free-map}, and the symbol map~\eqref{eq:uc-symbolmap} give the comparison.
Every chain arrow has its displayed source and target, including the unit-complex summand.
They also specify the next maps after the binary collision: external products on three factors, with compatibility supplied by the chiral Jacobi identity.
The two scalar collisions in~\eqref{eq:uc-firstcoll} lie below their input PBW orders and are retained in the filtered complex.
Passing to the symbol does not assert their vanishing before specialization.

For the fixed-coordinate relative state complex
\[
 C_{\mathrm{rel}}=\ker b_0\cap\ker H\subset A,
\]
the four-state calculation gives $C_{\mathrm{rel}}=\mathbb Ce\oplus\mathbb Cv$ with zero differential, in degrees zero and two.
It lies in PBW order zero.
Consequently it cannot be quasi-isomorphic, with these degrees, to $C^{\mathrm{ch}}(\mathbb P^1,\mathcal A)$ or to its PBW associated graded complex.
Their cohomology degrees are those of Theorem~\ref{thm:uc-global}.
This comparison concerns the unmarked proper curve and algebraic, uncompleted chiral chains.
Marked-point coefficients, relative geometric constructions, and inverse-limit completions require their own carriers and comparison maps.
